\section{Dataset}
The dataset we are working with is the \href{https://citibikenyc.com/system-data}{Citi Bike System Data}, which contains information about bike rentals in New York City.
The data are provided by the private company \textit{Lyft}, which operates the bike-sharing service in the city.
\\
All data are collected and made open source by the company, updated on a monthly basis, and made available through an API.
This means that we have access to a continuous stream of data, which allows us to analyze trends and patterns over time.
\\
The dataset contains information such as:
\begin{itemize}
    \item Rental start and end times
    \item Rental start and end stations
    \item User type (e.g., subscriber or customer)
\end{itemize}
The data are available for the last 10 years, and they are provided in a structured format (CSV files), which makes them easier to process and analyze.

\subsection{Other Datasets}
Since weather is a factor that can influence mobility patterns, we also plan to integrate weather data into our analysis.
This is possible through the use of APIs of \href{https://openweathermap.org/api}{OpenWeatherMap}, which can provide historical weather data for New York City.
By combining the bike rental data with weather data, we can gain a deeper understanding of how weather conditions affect mobility patterns and bike usage in the city.
\\\\
We also take into consideration \href{https://www.nyc.gov/html/dot/html/about/datafeeds.shtml\#Bikes}{NYC Bicycling Data}, which can provide additional information such as bike lane locations, traffic data, and other relevant context that can help us understand bike rentals and mobility patterns in New York City.
We can also consider other datasets that provide information about the city's infrastructure, such as public transportation data, to gain a more comprehensive understanding of mobility patterns.